\documentclass[12pt,a4paper]{report}

% ======= Configuración básica =======
\usepackage[spanish, es-tabla]{babel}
\usepackage[utf8]{inputenc}
\usepackage[T1]{fontenc}
\usepackage{lmodern}
\usepackage{geometry}
% Márgenes ajustados (tócalos si quieres aún más estrecho/ancho)
\geometry{
  a4paper,
  left=25mm,
  right=18mm,
  top=20mm,
  bottom=22mm,
  includeheadfoot
}
\usepackage{setspace}
\onehalfspacing
\usepackage{graphicx}
\graphicspath{{figs/}}
\usepackage{booktabs}
\usepackage{amsmath,amssymb}
\usepackage{csquotes}
\usepackage[hidelinks]{hyperref}
\usepackage[nottoc]{tocbibind} % Añade LoF al TOC

% ======= Metadatos =======
\newcommand{\titulo}{BasketballAnalysis}
\newcommand{\subtitulo}{Big Data de la NBA con IA y modelos predictivos de \textit{machine learning}}
\newcommand{\autor}{Pablo Díaz}
\newcommand{\fecha}{\today}

% Opcional: pequeñas mejoras tipográficas
\setlength{\parskip}{0.35em}
\setlength{\parindent}{12pt}

\begin{document}

% ========= Portada =========
\begin{titlepage}
  \centering
  {\huge \bfseries \titulo \par}
  \vspace{0.6cm}
  {\large \subtitulo \par}
  \vfill
  {\Large \autor \par}
  \vspace{0.4cm}
  {\large \fecha \par}
\end{titlepage}

% ========= Preliminares =========
\pagenumbering{Roman}
\tableofcontents
\cleardoublepage
\listoffigures
\cleardoublepage

% ========= Acrónimos (tu capítulo en itemize) =========
\cleardoublepage
\phantomsection
\chapter*{Acrónimos}
\addcontentsline{toc}{chapter}{Acrónimos}

\begin{itemize}
  % Generales
  \item \textbf{NBA} — National Basketball Association
  \item \textbf{IA} — Inteligencia Artificial
  \item \textbf{ML} — \textit{Machine Learning}
  \item \textbf{EDA} — \textit{Exploratory Data Analysis}
  \item \textbf{API} — \textit{Application Programming Interface}
  \item \textbf{ETL} — Extract, Transform, Load

  % Métricas de baloncesto
  \item \textbf{eFG\%} — \textit{Effective Field Goal Percentage}
  \item \textbf{TS\%} — \textit{True Shooting Percentage}
  \item \textbf{ORtg} — \textit{Offensive Rating}
  \item \textbf{DRtg} — \textit{Defensive Rating}
  \item \textbf{NetRtg} — Net Rating
  \item \textbf{PACE} — Ritmo (posesiones por 48 min)
  \item \textbf{RAPM} — \textit{Regularized Adjusted Plus-Minus}

  % Modelado
  \item \textbf{RF} — \textit{Random Forest}
  \item \textbf{XGBoost} — \textit{Extreme Gradient Boosting}
  \item \textbf{ROC} — \textit{Receiver Operating Characteristic}
  \item \textbf{AUC} — Área bajo la curva ROC
  \item \textbf{RMSE} — Raíz del Error Cuadrático Medio
  \item \textbf{MAE} — Error Absoluto Medio
  \item \textbf{SHAP} — \textit{SHapley Additive exPlanations}

  % Proyecto
  \item \textbf{W/L} — \textit{Win/Loss}
\end{itemize}


% ========= Cuerpo =========
\pagenumbering{arabic}
\chapter{Introducción}\label{chap:introduccion}

\textbf{BasketballAnalysis} es un proyecto orientado a aplicar técnicas modernas de \gls{ia} y \gls{ml} sobre grandes volúmenes de datos de la \gls{nba}. 
El objetivo es construir un flujo de trabajo reproducible que permita (i) entender qué factores explican el rendimiento de equipos, alineaciones y jugadores, 
y (ii) entrenar modelos predictivos capaces de estimar resultados \gls{wl} y métricas asociadas al desempeño.

\section*{Motivación}
La \gls{nba} ofrece un ecosistema rico en datos: estadísticas tradicionales y avanzadas, \textit{play-by-play}, información de contextos (local/visitante, descansos, parciales), e incluso datos derivados (on/off, alineaciones, clústeres de tiro). 
Aprovechar este \textit{Big Data} con \gls{ia}/\gls{ml} permite pasar de la descripción a la predicción, optimizando la toma de decisiones (rotaciones, emparejamientos, recursos) y generando conocimiento accionable.

\section*{Problema a resolver}
Dado un conjunto heterogéneo de fuentes (dashboards de equipo, alineaciones, tiro, pase, rebote, on/off, etc.), se requiere:
\begin{enumerate}
  \item Integrar, limpiar y estandarizar tablas para un modelo de datos coherente.
  \item Definir y calcular métricas clave (\gls{efg}, \gls{ts}, \gls{ortg}, \gls{drtg}, \gls{netrtg}, \gls{pace}) y derivadas contextuales (descanso, rachas, calendario).
  \item Identificar \textit{drivers} del rendimiento (correlaciones, importancia de variables, análisis de sensibilidad).
  \item Entrenar y evaluar modelos predictivos para resultados \gls{wl} y otros objetivos (p.~ej. diferencial de puntos).
\end{enumerate}

\section*{Objetivos específicos}
\begin{itemize}
  \item Diseñar un pipeline reproducible de \gls{etl} y \gls{eda} sobre datos de la \gls{nba}.
  \item Construir y validar modelos de clasificación (p.~ej. \gls{rf}, \gls{xgb}) para predicción \gls{wl}.
  \item Evaluar modelos con métricas robustas (\gls{auc}, \gls{mae}, \gls{rmse}, calibración) y análisis de importancia/explicabilidad (\gls{shap}).
  \item Documentar supuestos, limitaciones y recomendaciones de uso.
\end{itemize}

\section*{Alcance}
El documento cubre la cadena completa, desde la descripción de datos y métricas, hasta la preparación de \textit{features}, 
\textit{modeling} y evaluación. No se busca sustituir la \textit{scouting} cualitativa, sino complementarla con evidencia cuantitativa.

\section*{Resumen metodológico}
El flujo propuesto:
\begin{enumerate}
  \item \textbf{Ingesta \& limpieza}: carga de \texttt{parquet}, normalización de columnas y llaves (p.~ej. \texttt{GAME\_ID}, \texttt{TEAM\_ID}).
  \item \textbf{Ingeniería de variables}: métricas avanzadas, contexto (local/visitante, descansos, post-\textit{All-Star}), \textit{rolling windows}, \textit{on/off}, alineaciones.
  \item \textbf{EDA}: correlaciones (Spearman), mapas de calor, \textit{feature screening}.
  \item \textbf{Modelado}: \gls{rf}/\gls{xgb} con validación temporal, cuidado del \textit{leakage}, balanceo si procede.
  \item \textbf{Evaluación}: \gls{auc}, \gls{mae}/\gls{rmse} (si regresión), curva \gls{roc}, calibración y \gls{shap}.
\end{enumerate}

\section*{Organización del documento}
El Capítulo~\ref{chap:marco} presenta el marco teórico; el Capítulo~\ref{chap:datos} describe los datos; 
el Capítulo~\ref{chap:metodologia} detalla el pipeline de preparación y modelado; 
el Capítulo~\ref{chap:resultados} muestra resultados y visualizaciones; y el Capítulo~\ref{chap:conclusiones} cierra con conclusiones y líneas futuras.

\chapter{Marco teórico}\label{chap:marco}

Este capítulo establece los fundamentos conceptuales necesarios para el proyecto \textit{BasketballAnalysis}. 
Primero se describen conceptos clave del baloncesto cuantitativo (posesiones, eficiencia, cuatro factores, 
alineaciones y contexto competitivo). Después, se presenta el armazón de aprendizaje automático empleado, 
incluyendo formulación del problema, validación temporal, modelos elegidos y consideraciones de calibración e interpretabilidad.

\section{Baloncesto: fundamentos y métricas}\label{sec:baloncesto}

\subsection*{Posesiones y ritmo}
El análisis comparativo exige normalizar por volumen de posesiones. Una aproximación habitual al conteo de posesiones de un equipo es:
\begin{equation}
\mathrm{Poss} \approx \mathrm{FGA} + 0.44\,\mathrm{FTA} - \mathrm{OREB} + \mathrm{TOV}.
\end{equation}
El \textbf{ritmo} o \textbf{PACE} expresa posesiones por 48 minutos y permite comparar equipos que juegan a velocidades distintas.

\subsection*{Eficiencia de tiro}
La \textbf{eFG\%} ajusta el valor del triple:
\begin{equation}
\mathrm{eFG\%} = \frac{\mathrm{FGM} + 0.5 \cdot \mathrm{FG3M}}{\mathrm{FGA}}.
\end{equation}
La \textbf{TS\%} aproxima la eficiencia de anotación global incorporando tiros libres:
\begin{equation}
\mathrm{TS\%} = \frac{\mathrm{PTS}}{2\,(\mathrm{FGA} + 0.44\,\mathrm{FTA})}.
\end{equation}

\subsection*{Ratings por 100 posesiones}
\begin{equation}
\mathrm{ORtg} = 100 \cdot \frac{\mathrm{PTS}}{\mathrm{Poss}}, \qquad
\mathrm{DRtg} = 100 \cdot \frac{\mathrm{PTS\ concedidos}}{\mathrm{Poss}}, \qquad
\mathrm{NetRtg} = \mathrm{ORtg} - \mathrm{DRtg}.
\end{equation}
Estos indicadores permiten separar contribuciones ofensivas y defensivas normalizadas.

\subsection*{Cuatro factores}
Según la descomposición clásica, el rendimiento ofensivo se resume en cuatro pilares:
(i) eficacia de tiro (eFG\%), (ii) pérdidas (TOV\%), (iii) rebote ofensivo (OREB\%) y (iv) tiros libres (FTr).
Analizar su contribución relativa ayuda a explicar variaciones en \textit{NetRtg} entre equipos, alineaciones y contextos.

\subsection*{Alineaciones, on/off y estabilidad}
La composición de quintetos modula los espacios de tiro, la creación y la contención defensiva.
Las métricas \textit{on/off} estiman el impacto de un jugador comparando el rendimiento del equipo cuando está en pista
frente a cuando no lo está. No obstante, requieren \textit{shrinkage} o regularización por tamaño muestral para evitar conclusiones espurias.
La \textbf{estabilidad de rotación} (varianza de quintetos utilizados) es una pista útil de consistencia táctica.

\subsection*{Contexto competitivo}
El lugar del partido (local/visitante), los días de descanso, los \textit{back-to-back} y el calendario reciente condicionan el rendimiento.
Las segmentaciones temporales (meses, pre/post All-Star) capturan cambios de tendencia y ajustes de plantilla.

\section{Modelos y aprendizaje automático}\label{sec:modelos}

\subsection*{Formulación del problema}
El proyecto incluye estimadores de distinta naturaleza. El \textbf{Estimador 1} se plantea como clasificación binaria (W/L por partido).
Otros estimadores pueden formularse como regresión (diferencial de puntos) o agregación/proyección (victorias de temporada).

\subsection*{Representación de datos y prevención de \textit{leakage}}
Se construye una tabla de \textbf{partido–equipo}, donde cada fila representa un equipo en un partido dado.
Todas las variables (\textit{features}) se calculan usando únicamente información disponible \emph{antes} del salto inicial.
Se emplean diferencias simétricas frente al rival (\texttt{Feature\_Diff = Equipo - Rival}) para centrar el modelo en ventajas relativas.

\subsection*{Validación temporal}
La evaluación se realiza con \textbf{backtesting \textit{walk-forward}}: se entrena en el histórico hasta un punto temporal $T$ y se valida en el intervalo $(T, T+\Delta]$.
Este esquema respeta la causalidad y evita fugas de información propias de particiones aleatorias.

\subsection*{Modelos utilizados}
\paragraph{Random Forest (RF).}
Bosque de árboles con \textit{bagging}. Ventajas: robustez, poco \textit{tuning} y buena línea base.
Parámetros prácticos: número de árboles alto, profundidad moderada, tamaño mínimo de hoja suficiente para evitar sobreajuste.

\paragraph{XGBoost (XGB).}
\textit{Gradient boosting} con regularización explícita.
Capta no linealidades e interacciones de forma eficiente y suele optimizar métricas como LogLoss con buen sesgo-varianza.
Parámetros clave: \texttt{learning\_rate}, \texttt{max\_depth}, \texttt{subsample}, \texttt{colsample\_bytree}, \texttt{min\_child\_weight}, y regularización L1/L2.

\subsection*{Ponderación temporal y calibración}
Las observaciones más recientes se ponderan más (decaimiento exponencial por días).
La salida probabilística se \textbf{calibra} (Platt o isotónica) para que los porcentajes predichos coincidan con frecuencias observadas,
lo cual es crucial en toma de decisiones.

\subsection*{Interpretabilidad}
Se emplean importancias por permutación y, en su caso, \textit{SHAP values} para explicar contribuciones globales y locales.
La interpretabilidad permite auditar el comportamiento del modelo y detectar variables espurias o inestables.

\subsection*{Despliegue y control de calidad}
En producción, el flujo actualiza las \textit{features} hasta la víspera, infiere para los partidos del día y genera reportes de rendimiento,
calibración y deriva (cambios de distribución) para mantenimiento preventivo.

\chapter{Datos}\label{chap:datos}

En este capítulo se describe el origen de los datos, esquema de campos y controles de calidad (duplicados, rangos, nulos).
% Ejemplo de figura (opcional):
% \begin{figure}[h]
%   \centering
%   \includegraphics[width=.8\textwidth]{../figs/ejemplo.pdf}
%   \caption{Diagrama de flujo del dataset.}
%   \label{fig:flujo-datos}
% \end{figure}

\chapter{Metodología}

Se detallan pasos de limpieza, reglas de imputación, cálculo de métricas avanzadas y criterios para comparar lineups y jugadores.

\chapter{Modelos predictivos (Estimadores)}\label{chap:estimadores}

Este capítulo describe la familia de estimadores que convivirán en el proyecto. 
El \textbf{Estimador 1} aborda la probabilidad de victoria por partido (W/L).
Se dejan además dos plantillas extensibles para estimadores adicionales, de modo que el proyecto crezca de forma coherente y modular.

\section{Visión general}\label{sec:vision-general}
Todos los estimadores comparten un \textbf{núcleo de datos} y criterios de calidad:
(i) ingesta y normalización de fuentes (\textit{dashboards}, alineaciones, on/off, tiro/pase/rebote),
(ii) ingeniería de variables por bloques, (iii) prevención de \textit{leakage}, 
(iv) validación temporal \textit{walk-forward}, (v) calibración de salidas y (vi) reportes con explicabilidad.

\section{Estimador 1: Victoria/Derrota por partido (W/L)}\label{sec:estimador1}

\subsection{Objetivo y unidad de análisis}
La unidad es \textbf{partido–equipo}. Para cada encuentro se generan dos registros (local y visitante).
La variable objetivo es binaria: $y=1$ si el equipo gana y $y=0$ si pierde.
El modelo produce una probabilidad \texttt{p\_win} previa al salto inicial y, opcionalmente, una proyección de victorias a futuro.

\subsection{Diseño de \textit{features}: bloques y filosofía}
Las \textit{features} se organizan en bloques, cada una calculada \emph{hasta la víspera} del partido y en variantes \textit{rolling} (3/5/10 últimos partidos) y \textit{season-to-date}. 
Se prioriza la \textbf{simetría} mediante diferencias frente al rival (\texttt{Equipo - Rival}) para capturar ventajas relativas.

\paragraph{Bloque A — Fortaleza base y forma}
\textit{Off/Def/Net Rating} y \textit{PACE}, junto con los \textbf{Cuatro Factores} (eFG\%, TOV\%, OREB\%, FTr). 
Se incluyen tendencias (diferencias recientes o pendientes) para reflejar forma.

\paragraph{Bloque B — Perfil de tiro y contestación}
Distribuciones por zona (aro, pintura, media, esquinas, \textit{above-the-break}), 
tipología (pull-up vs. catch\&shoot) y contestación (distancia del defensor, tiempo de toque). 
Se construyen señales de \textbf{calidad esperada} (xEFG) y se contrasta con lo que el rival concede.

\paragraph{Bloque C — Rebote y segundas oportunidades}
\textit{OREB\%}, \textit{DREB\%}, cadenas de segundas oportunidades y, si está disponible, métricas de \textit{hustle} (boxouts, balones sueltos).

\paragraph{Bloque D — Creación y pérdidas}
\textit{Assist\%}, asistencias potenciales y secundarias, pases por posesión, puntos creados por pase, 
\textit{TOV\%} (idealmente separando pérdidas en vivo/muertas), transición (\% posesiones y puntos).

\paragraph{Bloque E — Tiros libres y disciplina}
\textit{Free Throw Rate} a favor/en contra, \textit{FT\%} y balance de faltas (recibidas/cometidas).

\paragraph{Bloque F — Contexto y calendario}
Local/visitante, diferencia de descanso, \textit{back-to-back}, tramos de 3 en 4 días y segmentaciones temporales (meses, pre/post All-Star).

\paragraph{Bloque G — Clutch con \textit{shrinkage}}
\textit{NetRtg}, eFG\% y pérdidas en minutos clutch, regularizados fuertemente para evitar ruido por muestras pequeñas.

\paragraph{Bloque H — Alineaciones esperadas y disponibilidad}
Quinteto probable y rotación corta. 
\textit{Net/Off/Def Rating} de esas combinaciones (season-to-date y \textit{rolling}) ponderados por minutos.
\textit{On/Off} de piezas clave con penalización por tamaño muestral.
Estabilidad de rotación (varianza de quintetos utilizados recientemente).

\paragraph{Bloque I — Matchups estilísticos}
Compatibilidad entre lo que busca el ataque propio y lo que permite la defensa rival (rim attempts, corner 3, PnR handler, poste, etc. si están disponibles).

\subsection{Construcción del dataset y control de calidad}
\begin{enumerate}
  \item Generar el calendario de partidos con identificadores (\texttt{GAME\_ID}, fecha, local/visitante, equipos).
  \item Para cada fecha, calcular todos los bloques hasta la víspera y cruzar contra el rival para obtener diferencias.
  \item Etiquetar el resultado W/L para cada registro partido–equipo.
  \item Imputar valores ausentes con estadísticas robustas (medianas por equipo o globales) y winsorizar \textit{outliers}.
  \item Asignar pesos temporales $w=\exp(-\Delta\text{días}/\text{halflife})$ (30–45 días como guía).
\end{enumerate}

\subsection{Validación y selección de modelos}
El rendimiento se evalúa con \textbf{backtesting \textit{walk-forward}} (entrenar hasta $T$, validar en $(T, T+\Delta]$).
Métricas principales: \textit{LogLoss}, \textit{Brier score} y AUC, además de análisis de calibración (curvas de fiabilidad, ECE).

\paragraph{RF (línea base robusta).}
Árboles numerosos, profundidad moderada, hojas con tamaño mínimo generoso. 
Importancias por permutación en las ventanas de validación.

\paragraph{XGB (potencia y control).}
\textit{Learning rate} bajo, \textit{subsample} y \textit{colsample} en rangos conservadores, \textit{early stopping}.
Regularización L1/L2 y \texttt{min\_child\_weight} para limitar complejidad efectiva.

\paragraph{Ensemble y calibración.}
Combinación ponderada de RF y XGB (según \textit{LogLoss} en validación) o \textit{stacking} con una logística final entrenada sobre predicciones \textit{out-of-fold}.
Calibración posterior (Platt o isotónica) en el último tramo de validación.

\subsection{Producción y reportes operativos}
\begin{itemize}
  \item Actualizar diariamente las \textit{features} hasta la víspera y obtener \texttt{p\_win} de los partidos del día.
  \item Proyección de temporada: suma de probabilidades o simulación Monte Carlo sobre el calendario restante.
  \item Paneles de control: estabilidad mensual, sesgos local/visitante, bins de probabilidad vs. frecuencia real (calibración).
  \item Auditoría: importancias por permutación y \textit{SHAP} (si se usa XGB) a nivel global y por partido.
\end{itemize}

\subsection{Checklist mínimo viable (MVP)}
\begin{itemize}
  \item Núcleo: \texttt{NetRtg\_diff}, \texttt{eFG\%\_diff}, \texttt{TOV\%\_diff}, \texttt{OREB\%\_diff}, \texttt{FTr\_diff}.
  \item Apoyo: \texttt{3PA rate diff}, \texttt{rim attempt diff}, \texttt{opp rim FG\% allowed diff}.
  \item Contexto: \texttt{Rest\_diff}, \texttt{home}, flags de \textit{back-to-back}.
  \item Plus: \texttt{Clutch\_NetRtg\_diff} regularizado y \texttt{Lineup\_expected\_NetRtg\_diff}.
\end{itemize}

\section{Estimador 2: Diferencial de puntos por partido}\label{sec:estimador2}
\textbf{Objetivo (regresión).} Predecir el diferencial de puntos final (\texttt{pts\_for} - \texttt{pts\_against}) en la unidad partido–equipo. 
Comparte bloques A–I del Estimador 1, pero el \textbf{target} es continuo. 
Se evalúa con RMSE/MAE y se vigila la calibración \textit{quantile} (residuos simétricos).

\textbf{Modelos.} RF regresión (línea base) y XGB regresión (ajustando \texttt{max\_depth}, \texttt{learning\_rate}, regularización).
Atención al \textbf{sesgo local/visitante} y a la \textbf{heterocedasticidad}: 
puede añadirse un modelo de varianza condicional o transformar el target (p.ej., winsorizar).

\textbf{Uso.} Este estimador alimenta pronósticos derivados (márgenes, \textit{spreads}) y explica de forma continua el grado de superioridad/inferioridad esperado.

\section{Estimador 3: Proyección de victorias de temporada}\label{sec:estimador3}
\textbf{Objetivo (agregación/simulación).} Convertir probabilidades partido a partido en \textbf{victorias esperadas} o 
en distribuciones de victorias mediante \textbf{Monte Carlo}. 

\textbf{Flujo.} Usar las \texttt{p\_win} del Estimador 1 a lo largo del calendario restante; 
(1) sumar probabilidades para una expectativa puntual, y/o (2) simular $N$ trayectorias de temporada para obtener percentiles (P10/P50/P90).

\textbf{Control.} Recalibrar periódicamente \texttt{p\_win} con resultados recientes, 
monitorizar deriva por cambios de plantilla y documentar supuestos (lesiones, descansos, traspasos).

\section{Cierre}
Este conjunto de estimadores cubre predicción binaria (W/L), magnitud del resultado (diferencial) y objetivos de planificación (victorias de temporada).
La modularidad facilita añadir variantes (p.ej., \textit{in-game live}, métricas por alineación, o \textit{player impact}) sin reescribir el armazón.
 % <--- NUEVO capítulo extenso
\chapter{Conclusiones y trabajo futuro}

Resumen de aprendizajes clave, limitaciones y siguientes pasos (mejoras de datos, nuevos modelos, dashboards, etc.).


\end{document}
